% GENERATED FILE. Edit canonical lessons, then run npm run generate:books.
% canonical-source-hash: fnv1a64:fcdd203ddc66ff81
% canonical-lessons: TE-C432-niti, TE-C432-adarsam, TE-C432-pratistha, TE-C432-visvasaniyata, TE-C432-avagahana

\chapter{Ideals and Standing}
\label{ch:te-ideals-and-standing}
\hlchaptermodality{432}

\begin{chapteropening}
\textbf{By the end of this chapter:} \emph{I can talk about ethics, ideals, reputation and trust.}

Complete the last lesson of chapter 432: I can talk about ethics, ideals, reputation and trust.
\end{chapteropening}

\section[\texorpdfstring{\te{నీతి} (nīti)}{నీతి (nīti)}]{\te{నీతి} (nīti) — ethics; a moral}

\label{lesson:TE-C432-niti}



\begin{quote}
Before the new one: say the Telugu for effective, then the Telugu for to test out.
\end{quote}

\subsection*{You'll want to know: \te{నీతి}}
\textbf{\te{నీతి}} — \emph{nīti} — \textquotedblleft{}ethics; a moral\textquotedblright{}.

\begin{quote}
\te{ఆ} \te{కథలో} \te{మంచి} \te{నీతి} \te{ఉంది}.
\end{quote}

\begin{quote}
\emph{ā kathalō mañci nīti undi.}
\end{quote}

\begin{quote}
That story has a good moral.
\end{quote}

\te{నీతి} \te{కథ}, a fable with a moral; \te{నీతిగా}, honestly.

\begin{cousinweb}
\textbf{\te{నీతి}} (\emph{nīti}) is borrowed from Sanskrit \emph{nīti}, \textquotedblleft{}conduct\textquotedblright{} (CDIAL 7546), and especially \textquotedblleft{}right or wise conduct\textquotedblright{} (Monier-Williams).
\end{cousinweb}

\begin{grammarlens}[title={What you've built}]
One more word for saying what happens if.
\end{grammarlens}

\subsection*{Your turn}
\begin{itemize}
\raggedright
  \item \emph{Say it:} \emph{nīti}
  \item \emph{Say it:} \emph{nīti}, once more
  \item \emph{Say it:} \emph{ā kathalō mañci nīti undi.} — That story has a good moral.
  \item \emph{Recall:} say the Telugu for effective, then the Telugu for to test out, then say \emph{nīti} again
\end{itemize}

\begin{tcolorbox}[breakable,colback=teal!4,colframe=teal!35!black,title={Before you move on}]
What does \te{నీతి} mean? (\textquotedblleft{}Ethics; a moral\textquotedblright{}.) Say it once more.
\end{tcolorbox}
\section[\texorpdfstring{\te{ఆదర్శం} (ādarśaṁ)}{ఆదర్శం (ādarśaṁ)}]{\te{ఆదర్శం} (ādarśaṁ) — an ideal; a model}

\label{lesson:TE-C432-adarsam}



\begin{quote}
Before the new one: say the Telugu for to test out, then the Telugu for ethics.
\end{quote}

\subsection*{You'll want to know: \te{ఆదర్శం}}
\textbf{\te{ఆదర్శం}} — \emph{ādarśaṁ} — \textquotedblleft{}an ideal; a model\textquotedblright{}.

\begin{quote}
\te{మా} \te{అమ్మే} \te{నాకు} \te{ఆదర్శం}.
\end{quote}

\begin{quote}
\emph{mā ammē nāku ādarśaṁ.}
\end{quote}

\begin{quote}
My mother is my role model.
\end{quote}

Before a noun, \te{ఆదర్శ}: \te{ఆదర్శ} \te{గ్రామం}, a model village.

\begin{cousinweb}
\textbf{\te{ఆదర్శం}} (\emph{ādarśaṁ}) is borrowed from Sanskrit \emph{ādarśa}, \textquotedblleft{}a mirror\textquotedblright{} (CDIAL 1143); Telugu uses it for a model to follow.
\end{cousinweb}

\begin{grammarlens}[title={What you've built}]
One more word for saying what happens if.
\end{grammarlens}

\subsection*{Your turn}
\begin{itemize}
\raggedright
  \item \emph{Say it:} \emph{ādarśaṁ}
  \item \emph{Say it:} \emph{ādarśaṁ}, once more
  \item \emph{Say it:} \emph{mā ammē nāku ādarśaṁ.} — My mother is my role model.
  \item \emph{Recall:} say the Telugu for to test out, then the Telugu for ethics, then say \emph{ādarśaṁ} again
\end{itemize}

\begin{tcolorbox}[breakable,colback=teal!4,colframe=teal!35!black,title={Before you move on}]
What does \te{ఆదర్శం} mean? (\textquotedblleft{}An ideal; a model\textquotedblright{}.) Say it once more.
\end{tcolorbox}
\section[\texorpdfstring{\te{ప్రతిష్ఠ} (pratiṣṭha)}{ప్రతిష్ఠ (pratiṣṭha)}]{\te{ప్రతిష్ఠ} (pratiṣṭha) — prestige; reputation}

\label{lesson:TE-C432-pratistha}



\begin{quote}
Before the new one: say the Telugu for ethics, then the Telugu for an ideal.
\end{quote}

\subsection*{You'll want to know: \te{ప్రతిష్ఠ}}
\textbf{\te{ప్రతిష్ఠ}} — \emph{pratiṣṭha} — \textquotedblleft{}prestige; reputation\textquotedblright{}.

\begin{quote}
\te{ఆ} \te{పని} \te{వల్ల} \te{కుటుంబ} \te{ప్రతిష్ఠ} \te{పెరిగింది}.
\end{quote}

\begin{quote}
\emph{ā pani valla kuṭumba pratiṣṭha perigindi.}
\end{quote}

\begin{quote}
That deed raised the family’s standing.
\end{quote}

\te{ప్రతిష్ఠాత్మకమైన}, prestigious.

\begin{cousinweb}
\textbf{\te{ప్రతిష్ఠ}} (\emph{pratiṣṭha}) is borrowed from Sanskrit \emph{pratiṣṭhā}, \textquotedblleft{}foundation, prop\textquotedblright{} (CDIAL 8608), and \textquotedblleft{}standing still, stability\textquotedblright{} (Monier-Williams); Telugu uses it for standing in people’s eyes.
\end{cousinweb}

\begin{grammarlens}[title={What you've built}]
One more word for saying what happens if.
\end{grammarlens}

\subsection*{Your turn}
\begin{itemize}
\raggedright
  \item \emph{Say it:} \emph{pratiṣṭha}
  \item \emph{Say it:} \emph{pratiṣṭha}, once more
  \item \emph{Say it:} \emph{ā pani valla kuṭumba pratiṣṭha perigindi.} — That deed raised the family’s standing.
  \item \emph{Recall:} say the Telugu for ethics, then the Telugu for an ideal, then say \emph{pratiṣṭha} again
\end{itemize}

\begin{tcolorbox}[breakable,colback=teal!4,colframe=teal!35!black,title={Before you move on}]
What does \te{ప్రతిష్ఠ} mean? (\textquotedblleft{}Prestige; reputation\textquotedblright{}.) Say it once more.
\end{tcolorbox}
\section[\texorpdfstring{\te{విశ్వసనీయత} (viśvasanīyata)}{విశ్వసనీయత (viśvasanīyata)}]{\te{విశ్వసనీయత} (viśvasanīyata) — credibility; reliability}

\label{lesson:TE-C432-visvasaniyata}



\begin{quote}
Before the new one: say the Telugu for an ideal, then the Telugu for prestige.
\end{quote}

\subsection*{You'll want to know: \te{విశ్వసనీయత}}
\textbf{\te{విశ్వసనీయత}} — \emph{viśvasanīyata} — \textquotedblleft{}credibility; reliability\textquotedblright{}.

\begin{quote}
\te{ఆ} \te{పత్రిక} \te{విశ్వసనీయత} \te{తగ్గింది}.
\end{quote}

\begin{quote}
\emph{ā patrika viśvasanīyata taggindi.}
\end{quote}

\begin{quote}
That paper’s credibility has fallen.
\end{quote}

\te{విశ్వసనీయమైన}, reliable, trustworthy; of the same family as \te{విశ్వసించు}.

\begin{cousinweb}
\textbf{\te{విశ్వసనీయత}} (\emph{viśvasanīyata}) is built on Sanskrit \emph{viśvasanīya}, \textquotedblleft{}to be trusted or relied on, reliable, credible\textquotedblright{} (Monier-Williams), with the abstract ending -\te{త}.
\end{cousinweb}

\begin{grammarlens}[title={What you've built}]
One more word for saying what happens if.
\end{grammarlens}

\subsection*{Your turn}
\begin{itemize}
\raggedright
  \item \emph{Say it:} \emph{viśvasanīyata}
  \item \emph{Say it:} \emph{viśvasanīyata}, once more
  \item \emph{Say it:} \emph{ā patrika viśvasanīyata taggindi.} — That paper’s credibility has fallen.
  \item \emph{Recall:} say the Telugu for an ideal, then the Telugu for prestige, then say \emph{viśvasanīyata} again
\end{itemize}

\begin{tcolorbox}[breakable,colback=teal!4,colframe=teal!35!black,title={Before you move on}]
What does \te{విశ్వసనీయత} mean? (\textquotedblleft{}Credibility; reliability\textquotedblright{}.) Say it once more.
\end{tcolorbox}
\section[\texorpdfstring{\te{అవగాహన} (avagāhana)}{అవగాహన (avagāhana)}]{\te{అవగాహన} (avagāhana) — awareness; understanding}

\label{lesson:TE-C432-avagahana}



\begin{quote}
Before the new one: say the Telugu for prestige, then the Telugu for credibility.
\end{quote}

\subsection*{You'll want to know: \te{అవగాహన}}
\textbf{\te{అవగాహన}} — \emph{avagāhana} — \textquotedblleft{}awareness; understanding\textquotedblright{}.

\begin{quote}
\te{జనంలో} \te{ఆరోగ్యం} \te{మీద} \te{అవగాహన} \te{పెరిగింది}.
\end{quote}

\begin{quote}
\emph{janaṁlō ārōgyaṁ mīda avagāhana perigindi.}
\end{quote}

\begin{quote}
People’s awareness of health has grown.
\end{quote}

\te{అవగాహన} \te{కల్పించు}, to raise awareness.

\begin{cousinweb}
\textbf{\te{అవగాహన}} (\emph{avagāhana}) is borrowed from Sanskrit \emph{avagāhana}, \textquotedblleft{}immersion, bathing\textquotedblright{} (Monier-Williams); Telugu uses it for being steeped in a subject, understanding.
\end{cousinweb}

\begin{grammarlens}[title={What you've built}]
One more word for saying what happens if.
\end{grammarlens}

\subsection*{Your turn}
\begin{itemize}
\raggedright
  \item \emph{Say it:} \emph{avagāhana}
  \item \emph{Say it:} \emph{avagāhana}, once more
  \item \emph{Say it:} \emph{janaṁlō ārōgyaṁ mīda avagāhana perigindi.} — People’s awareness of health has grown.
  \item \emph{Recall:} say the Telugu for prestige, then the Telugu for credibility, then say \emph{avagāhana} again
\end{itemize}

\begin{tcolorbox}[breakable,colback=teal!4,colframe=teal!35!black,title={Before you move on}]
What does \te{అవగాహన} mean? (\textquotedblleft{}Awareness; understanding\textquotedblright{}.) Say it once more.
\end{tcolorbox}
